\section{Related Work}

\textbf{Video Generation for Decision Making.} Recent advances in video models have achieved unprecedented visual quality and physical fidelity for video synthesis~\citep{guo2023animatediff, yang2024cogvideox, videoworldsimulators2024, veo2, wan2025}. This has demonstrated promise in summarizing world dynamics through videos~\citep{yang2024video, Bruce2024GenieGI} and has inspired the application of video models to solving decision-making problems~\citep{alejandro2023viper, du2024learning, yang2023learning, McCarthy2024TowardsGR, liang2024dreamitate}. Prior works have utilized video generative models as reward functions~\citep{luo2024text, alejandro2023viper, huang2023diffusion}, dynamics models~\citep{yang2023learning, Bruce2024GenieGI, Valevski2024DiffusionMA}, and pixel-based planners~\citep{ko2024avdc, ajaj2023hip, du2024video, Zhou2024RoboDreamerLC}. As in UniPi~\citep{du2024video}, we employ video models to predict text-conditioned visual plans that depict future outcomes, which are subsequently translated into actions via inverse dynamics. While the performance of such visual planners may often be limited by their offline pretraining data, our approach allows iterative improvement by learning from online environment interactions.





\textbf{Self-Improving Generative Models.} Continuously improving by learning from self-produced cumulative experience is an essential capability of intelligent agents. Prior work has demonstrated the effectiveness of improving LLMs with their self-generated outputs~\citep{yu-etal-2024-teaching, Tian2024NEURIPS_Toward_Self_Improvement, Huang2022CONFERENCE_Large_Language_Models}, where the LLM can serve as its own reward function~\citep{Yuan2024ICML_Self_Rewarding_Language} for preference optimization or data synthesizer~\citep {Patel2024ARXIV_Large_Language_Models} for supervised finetuning. However, a similar self-improvement recipe for video generation models remains underexplored. Most relevant to our work, VideoAgent~\citep{Soni2024VideoAgentSV} refines video generation through self-conditioning consistency and feedback from a VLM, and collects the successful plan rollouts for finetuning. We instead base our improvement loop on self-adaptation, where we leverage internet-scale video priors to synthesize improved visual plans for tasks unseen during initial in-domain training. Furthermore, our approach can still achieve self-improvement even with an initial model trained on suboptimal data and a notable relaxation on filtering requirements for finetuning data.

\textbf{Reinforcement Learning Finetuning of Behavior Cloning Policies.}
Behavior cloning approaches such as Diffusion Policy~\citep{chi2023diffusionpolicy}, which implements a policy as a diffusion model trained on offline-collected experience, are a performant approach for decision-making.  There have been numerous approaches for finetuning pretrained diffusion policies with respect to online experience and rewards.  ResIP~\citep{ankile2025imitation} utilizes a frozen diffusion policy model to propose action predictions, and learns a policy on top using reinforcement learning that transforms it into a more accurate action to perform in the environment.  DPPO~\citep{ren2024diffusion} treats the sampling procedure of a diffusion policy as an internal Markov Decision Process, and explicitly finetunes the weights of the pretrained diffusion policy with respect to achieved rewards.  DSRL~\citep{wagenmaker2025steering} utilizes a frozen diffusion policy in a deterministic fashion to map noise samples to action samples, and learns a noise selector through reinforcement learning.  In this work, we show that SILVR is more sample efficient than reinforcement learning finetuning of behavior cloning policies, and can achieve faster iterative improvements with respect to online experience.
